\scanpage{442}
\mspage{9}

a de plus ceci : si $G\in\mathrm{Ob}\,\uline{\mathrm{Grp}}(S)$ est un fibré en groupoïdes sur $S$, alors $G$ est équivalent à un fibré en groupoïdes de la forme $f^*(H)$ sss \deleted{\unsure{il existe|il est}} \inserted{pour} tout $t\in T$, la restriction \inserted{$G_t$} de $G_{\deleted{\unsure{t|s}}}$ à $\deleted{\unsure{T}}S_t$ \unsure{soit|est} équivalente à un \inserted{fibré en} groupoïdes constant sur $\deleted{\unsure{T}}S_t$. (\uline{NB} les conditions a) ou b) impliquent déjà que $S_t$ est $1$-connexe, i.e. que son groupoïde fondamental est \unsure{équiv.|équiv} à $e$.)

\textcircled{3} Soit $Z$ un objet $1$-connexe dans $(\mathrm{Cat})$, \deleted{Pour} $G$ un fibré en groupoïdes sur $Z$. Pour que $G$ soit \inserted{$Z$-}équivalent à un groupoïde constant sur $Z$, il \unsure{faut|f.} et s. que pour \unsure{tte|toute} flèche $u$ de $Z$, \deleted{\unsure{m.s.|m.s}} le \deleted{foncteur} morphisme de groupoïdes \unsure{cor\textsuperscript{t}|corr\textsuperscript{t}} induit sur les fibres de $G$ soit une \inserted{$Z$-}équivalence.

% slanted

\margin{On dit alors que $G$ est \uline{essentiellement} \uline{\unsure{loc.|loco.} constant} sur $Z$.}

\uline{Corollaire} La condition \unsure{trad.|trad.le} de $2^{\circ}$) s'exprime ainsi : pour toute flèche \inserted{$u$ de $S$} ``verticale'' \deleted{pour $f$} \inserted{sur $T$} i.e. telle que $f(u)$ soit une flèche identique, le foncteur $u^*$ induit \unsure{par} $u$ sur les fibres de $G$ soit une équivalence de groupoïdes

% archive sheet number: 441
